% Lineare Gleichungen – bank/lineare-gleichungen/ (zusammenbau v0.4, Heft)
\einheitenkopf[t9]{Lineare Gleichungen}
\setcounter{aufgabe}{64}

% Anlauf (ohne Prüfkennung)
\begin{aufgabe}{Lösung prüfen – Anlauf}
\begin{teile}
\teil $x + 9 = 15$ – ist $6$ Lösung?
\teil $2x + 7 = 13$ – ist $3$ Lösung?
\end{teile}
\end{aufgabe}

\begin{aufgabe}{Lösung prüfen – Prüfungsaufgaben}
\begin{teile}
\teil Welche Gleichung hat die Lösung $-3$? \hfill (P10 2018 OS) \\ \kreuz{$2x - 6 = 0$} \\ \kreuz{$-4x + 12 = 0$} \\ \kreuz{$3x + 5 = -4$} \\ \kreuz{$6x + 20 = 1$}
\teil Welche Gleichung hat die Lösung $-4$? \hfill (P10 2018 OS) \\ \kreuz{$5x - 20 = 0$} \\ \kreuz{$4x + 1 = 15$} \\ \kreuz{$2x + 11 = 3$} \\ \kreuz{$6x + 9 = -2$}
\teil $4 - 3x = 2x - 31$ – welche Zahl ist die Lösung? \hfill (P10 2022 OS) \\ \kreuz{$-7$} \\ \kreuz{$3$} \\ \kreuz{$7$} \\ \kreuz{$9$}
\teil $9 - 4x = 2x - 27$ – welche Zahl ist die Lösung? \hfill (P10 2022 OS) \\ \kreuz{$-6$} \\ \kreuz{$2$} \\ \kreuz{$6$} \\ \kreuz{$10$}
\end{teile}
\end{aufgabe}

% Anlauf (ohne Prüfkennung)
\begin{aufgabe}{Umformen – Anlauf}
\begin{teile}
\teil $x + 11 = 18$ – was kommt hinter den Strich? $\mid$ \leerfeld
\end{teile}
\begin{gleichungsraster}[2]
\gl{x + 8 = 15} &
\gl{x + 12 = 5} \\
\end{gleichungsraster}
\end{aufgabe}

\begin{aufgabe}{Umformen – Prüfungsaufgaben}
\begin{teile}
\teil $3(x + 7) = 12$ – löse und mache die Probe. \hfill (P10 2020 OS)
\teil $4(x + 9) = 20$ – löse und mache die Probe. \hfill (P10 2020 OS)
\teil $5 \cdot (x + 6) + 4 = 4$ – löse und mache die Probe. \hfill (P10 2023 OS)
\teil $4 \cdot (x + 2) - 7 = -7$ – löse und mache die Probe. \hfill (P10 2023 OS)
\teil $4 \cdot (x + 2{,}5) = 0$ – löse und mache die Probe. \hfill (P10 2024 OS)
\teil $3 \cdot (x + 1{,}5) = 0$ – löse und mache die Probe. \hfill (P10 2024 OS)
\end{teile}
\end{aufgabe}

% Anlauf (ohne Prüfkennung)
\begin{aufgabe}{Aufstellen – Anlauf}
\begin{teile}
\teil „Eine Zahl plus 9 ergibt 23." – welche Gleichung passt? ($x$: die gesuchte Zahl) \\ \kreuz{$9x = 23$} \\ \kreuz{$x + 9 = 23$} \\ \kreuz{$x - 9 = 23$}
\teil Ich denke mir eine Zahl. Addiere ich 18, erhalte ich 45. Gleichung und Zahl? ($x$: die gesuchte Zahl) x = \leerfeld
\end{teile}
\end{aufgabe}

\begin{aufgabe}{Aufstellen – Prüfungsaufgaben}
\begin{teile}
\steil Ein Freibad hatte am Samstag 380 Gäste, der Eintritt kostet 5 € pro Person. Der Betreiber sagt: „Mit $x$ Gästen mehr hätten wir 2\,400 € eingenommen." Stelle eine Gleichung auf und berechne, wie viele Gäste mehr es hätten sein müssen. \hfill (P10 2020 OS) \\ x = \leerfeld
\steil Ein Kino hat am Freitag 240 Karten zu je 9 € verkauft. Die Leiterin sagt: „Mit $x$ Karten mehr hätten wir 2\,700 € eingenommen." Stelle eine Gleichung auf und berechne, wie viele Karten mehr es hätten sein müssen. \hfill (P10 2020 OS) \\ x = \leerfeld
\teil 500 Jugendliche wurden nach ihrem Lieblingssport gefragt. Tennis wurde viermal so oft genannt wie Volleyball. Berechne die Anzahlen für Tennis und Volleyball. \hfill (P10 2018 OS) \\

\sachtabelle{lc}{Sport & Anzahl}{Fußball & 140\\ Schwimmen & 90\\ Basketball & 70\\ Sonstiges & 40\\ Tennis & \leerzelle\\ Volleyball & \leerzelle}
Volleyball: \leerfeld , Tennis: \leerfeld
\teil 800 Gäste eines Festes wurden nach ihrem Lieblingsgetränk gefragt. Kaffee wurde fünfmal so oft genannt wie Kakao. Berechne die Anzahlen für Kaffee und Kakao. \hfill (P10 2018 OS) \\

\sachtabelle{lc}{Getränk & Anzahl}{Wasser & 150\\ Saft & 210\\ Tee & 90\\ Limo & 110\\ Kaffee & \leerzelle\\ Kakao & \leerzelle}
Kakao: \leerfeld , Kaffee: \leerfeld
\teil „Das Neunfache einer Zahl vermindert um 15 ist gleich 57." Gib die Gleichung an. \hfill (P10 2016 OS) ($x$: die gesuchte Zahl)
\teil „Das Sechsfache einer Zahl vermindert um 20 ist gleich 34." Gib die Gleichung an. \hfill (P10 2016 OS) ($x$: die gesuchte Zahl)
\teil „Das Dreifache einer Zahl vermindert um 7 ist gleich 29." Welche Gleichung passt? \hfill (P10 2021 OS) ($x$: die gesuchte Zahl) \\ \kreuz{$3x = -7 + 29$} \\ \kreuz{$3 - 7x = 29$} \\ \kreuz{$3x - 7 = 29$} \\ \kreuz{$29 - 7 = 3x$}
\teil „Das Vierfache einer Zahl vermindert um 9 ist gleich 23." Welche Gleichung passt? \hfill (P10 2021 OS) ($x$: die gesuchte Zahl) \\ \kreuz{$4 - 9x = 23$} \\ \kreuz{$4x - 9 = 23$} \\ \kreuz{$23 - 9 = 4x$} \\ \kreuz{$4x = -9 + 23$}
\end{teile}
\end{aufgabe}
